\chapter{Methodology and Implementation}
\label{chap:methodology}

\Cref{chap:architecture} fixed what each component is responsible for and
what data it exchanges with its neighbours. This chapter explains how each
component was actually built: the concrete rules and libraries used, the
decisions made while building it, and -- where relevant -- the point at
which the real implementation forced a change to the design fixed in
\Cref{chap:architecture}. The work is organised into implementation steps
I1--I5, matching the project's own staged plan: I1 prepares the dataset,
I2 classifies and extracts context, I3 generates explanations
(building on the L4 model and prompt design), I4 generates repair proposals
(building on the L5 retrieval design), and I5 applies a proposal and
validates it by re-execution.

Two steps beyond I5 -- full pipeline orchestration and knowledge-graph
enrichment -- are not covered in this chapter because they have not been
built yet; they are addressed directly in \Cref{chap:future-work}.

\section{Dataset Preparation (I1)}
\label{sec:m-dataset}

\Cref{chap:problem-dataset} already gave the resulting numbers for step I1:
214 dependency-error rows, 200 usable and 14 excluded, split 13/187 between
development and evaluation. This section explains how
\texttt{scripts/prepare\_dependency\_dataset.py} produces those numbers, and
why the rules behind them were chosen.

The script reads the Docker pipeline's own \texttt{notebook\_executions}
table and keeps every row where \texttt{error\_category = 'DEPENDENCY\_ERROR'}.
For each kept row, it applies two independent decisions, deterministically
and without any LLM involvement: a \emph{subtype} (is the module missing
entirely, present but incompatible, a missing system library, or an
ambiguous local name?) and a \emph{scope status} (is this row inside the
pip-only v1 repair scope at all?). Both decisions rely only on simple,
inspectable rules over the \texttt{error\_type} and \texttt{error\_message}
fields the Docker pipeline already recorded -- for example, a message
containing \texttt{".so"} or "shared object file" is classified as
\texttt{system\_library}, and a failing module name drawn from a short list
of known-ambiguous names (\texttt{utils}, \texttt{statistics},
\texttt{config}, \texttt{helpers}, \texttt{src}) is classified as
\texttt{mapping\_unknown}.

Doing this filtering with plain rules, before any language model is
involved, was a deliberate choice. A decision that determines whether a row
is even eligible for automated repair should be explainable and
reproducible on its own terms; an LLM's judgement about scope eligibility
would be harder to audit and could silently change between runs of the same
model. Keeping this gate deterministic also means every later component can
trust \texttt{scope\_status} directly instead of re-deriving eligibility
itself.

The dataset is also split into three disjoint groups: 13 rows set aside for
manually designing and sanity-checking the explanation and repair prompts
(\Cref{sec:m-explanation,sec:m-repair}), 187 usable rows reserved for a
later, controlled evaluation run, and the 14 excluded rows. Separating a
small development sample from the evaluation set matters for the same
reason it matters in any machine-learning setting: a prompt tuned by looking
at a specific row, and then tested again on that same row, would give an
optimistic and misleading impression of how well the prompt generalises.

The script writes four output files: a CSV of the classified rows, a schema
document describing every column, a statistics file with the counts used
throughout this thesis, and a split-manifest file recording exactly which
notebook execution IDs fall into which split, so the same split can be
reproduced later rather than re-randomised.

\section{Context Extraction and Error Classification (I2)}
\label{sec:m-classification}

Step I2 (\texttt{scripts/extract\_error\_contexts.py}) takes each I1 row and
adds two things: a more detailed root-cause hint with a confidence level,
and, where possible, the actual source code around the failure. Every field
I1 already computed (\texttt{scope\_status}, \texttt{exclusion\_reason},
\texttt{split}) is carried through unchanged; I2 enriches every row
regardless of scope, so an excluded row still receives full context for its
explanation.

\subsection{Root-Cause Hints}

The subtype from I1 is refined into one of seven root-cause hints:

\begin{itemize}
  \item \texttt{direct\_missing\_package} -- the module is genuinely absent;
  \item \texttt{import\_distribution\_name\_mismatch} -- the import name
        differs from its PyPI distribution name, for example
        \texttt{sklearn} versus \texttt{scikit-learn};
  \item \texttt{version\_or\_api\_incompatibility} -- a "cannot import name"
        style failure;
  \item \texttt{transitive\_dependency} -- a generic \texttt{ImportError}
        not matching a more specific pattern;
  \item \texttt{system\_level\_dependency} -- a missing shared library;
  \item \texttt{local\_import\_or\_path\_issue} -- an ambiguous local module
        name; and
  \item \texttt{insufficient\_context} -- not enough information to decide.
\end{itemize}

Each hint is also assigned a confidence level -- high for a direct pattern
match (a missing-module message, a "cannot import name" message, or a
\texttt{.so} failure), medium for softer phrasing such as an "install" or
"required" hint, and low for an ambiguous local module or a message that
does not carry enough information to decide.

\subsection{Recovering Source Context}

Three tiers of context are attempted, in order of decreasing reliability.
The first checks a legacy execution table inherited from an earlier run of
the pipeline; because that table has no timestamp and was not produced by
the same execution that generated the current dataset, it is treated as a
hint, never as authoritative. A candidate row from that table is only
accepted once its own final reported error is verified to match the current
row -- not merely because the failing module's name happens to appear
somewhere in the cell source. This distinction matters in practice: a cell
that imports \texttt{pandas} before \texttt{sklearn} might fail on
\texttt{pandas}, while the row being classified records an \texttt{sklearn}
failure for the same notebook, and a naive substring match would wrongly
treat the two as the same failure. Cross-checking against this legacy table
shows that 162 of the 214 rows (76\%) have a matching notebook ID at all,
but only around 40\% of those actually mention the correct failing module in
their recorded traceback -- illustrating why this tier is used only as a
low-confidence hint and never trusted on its own.

The second tier, used only when explicitly requested, fetches the notebook
and its dependency files directly from GitHub, using the repository's
recorded commit where available and falling back to the \texttt{main} or
\texttt{master} branch otherwise. The third tier falls back to the metadata
already present from I1 (error type, message, cell index, failing module)
when no source can be recovered at all -- which is unavoidable here, because
the repositories cloned during the original Docker pipeline run were never
kept on this machine, so no local copy of any notebook's source exists
unless it is fetched again.

\subsection{Limitations}

A few limitations follow directly from the above and are worth stating
plainly, because they affect how confidently later components can rely on
this context. Remote fetching is not pinned to the exact commit that
originally failed unless that commit happens to still be resolvable, so a
freshly fetched notebook can, in rare cases, differ slightly from the one
that actually produced the recorded error. \texttt{error\_cell\_index}
indexes into the notebook's full cell list, so a notebook edited since the
original run can shift cell positions and misalign the index. Any
unexpected per-row failure (a malformed row, a database error) falls back
to a bare metadata-only record rather than aborting the whole batch, keeping
one bad row from blocking the other 213.

\section{LLM Explanation Component (I3)}
\label{sec:m-explanation}

\subsection{Model and Runtime Selection (L4)}

The explanation and repair components share one open, locally-run model
rather than each using a different one. \textbf{Gemma-2 9B} was selected as
the primary model because PLLM~\cite{bartlett2025pllm}, the closest prior
system, reports Gemma-2 9B with RAG as its own best-performing configuration
for exactly this kind of dependency task. Using the same model for both
explanation and repair keeps a later comparison between configurations
clean: any difference in outcome can be attributed to the prompt or the
retrieval design, not to a hidden difference in which model produced which
part of the result. CodeLlama-13B and DeepSeek-Coder (7B and 13B) were kept
as fallback candidates, motivated by Szalontai et
al.~\cite{szalontai2024codellama} and a broader empirical comparison of
code-specialised models~\cite{empiricalapr2025}, both of which report that
smaller, code-specialised models can match or beat larger general-purpose
ones on repair tasks; they were not needed in practice, since Gemma-2 9B
performed adequately in the design-time sanity checks described below.

Ollama was chosen as the serving runtime because it gives simple local
access to a quantised model without setting up a GPU-serving stack, which
is appropriate for a component running one notebook at a time rather than
serving high-throughput traffic. A hardware check confirmed \texttt{gemma2:9b}
runs on the available machine using 6.3~GB of memory, entirely on CPU, with
a 4096-token context window -- workable, though slow, which later became
relevant to a timeout decision described in \Cref{sec:m-repair}.

\subsection{From Prose to Structured JSON: an L4-to-I3 Deviation}
\label{sec:m-explanation-deviation}

The explanation prompt was originally designed (level L4) to produce
plain, unstructured prose: two to four short sentences, explicitly
forbidding headings, bullet points, or JSON. This was reasonable for a
first design-time sanity check, but it does not fit an automated pipeline
component whose output must be logged, validated, and machine-processed by
later steps. Once step I3 came to implement
\texttt{run\_llm\_explainer.py}, this requirement forced a change: the
explainer now returns a single JSON object, validated against a fixed
schema (\texttt{schemas/}\allowbreak\texttt{explanation.schema.json}) with
exactly six fields --
\texttt{summary}, \texttt{root\_cause}, an \texttt{evidence} list,
\texttt{failing\_module}, \texttt{explanation\_confidence} (high, medium, or
low), and \texttt{limitations} -- and no others.

This is a deviation from the original design, not a silent one: the L4
prose template remains valid evidence of the original design intent, and
the change was made because it was necessary, not because the original idea
was wrong. It reuses a pattern already validated for the repair prompt
(\Cref{sec:m-repair}): schema-constrained model output, followed by local
parsing and validation, with one retry on failure. The schema's
\texttt{explanation\_confidence} field is deliberately named differently
from I2's own classifier \texttt{confidence} field, because the two measure
different things and can legitimately disagree -- a classifier can be highly
confident it matched a known pattern while the model itself is only
moderately confident it has correctly explained the underlying cause.

\subsection{Prompt, Validation, and Retry}

The prompt instructs the model to write in plain, non-technical language,
to use only the information supplied (never inventing a version number, a
file name, or a cause not visible in the error), and never to suggest a fix
-- that is the repair component's job, kept structurally separate. The
runner (\texttt{run\_llm\_explainer.py}) sends this prompt to \texttt{gemma2:9b}
through Ollama, parses the response as JSON, and validates it against the
schema. On invalid JSON, a schema violation, a timeout, or a model
unavailability error, it retries once with a corrective prompt; if the
second attempt also fails, the failure itself is logged rather than
discarded, so that no row silently disappears from the results file. Each
logged result keeps three things separate: the original I2 input metadata,
the model and prompt configuration used, and the generated explanation
together with its validation status, latency, and retry count -- so that a
later reader can never confuse the original error text with text the model
generated about it.

\section{PyPI-Grounded Repair Generation (I4)}
\label{sec:m-repair}

RAGRepairAgent (\Cref{subsec:repairagent}) does not install anything. Its
only job is to answer one question -- \emph{what repair should be
attempted?} -- and to answer it safely: never inventing a package name,
never inventing a version number, and never handing a later component a raw
string to execute in a shell. This section explains how that is achieved.

\subsection{Overall Flow}

\begin{verbatim}
classified dependency failure (usable, from I2)
    -> identify the failing import
    -> resolve the import name to a PyPI distribution name
    -> query the PyPI registry for that distribution
    -> filter the returned versions to a small, safe candidate set
    -> for a version-incompatibility failure, intersect with
       curated compatibility evidence
    -> give the LLM only this grounded candidate set
    -> validate the LLM's proposal against the same evidence
    -> if valid, deterministically build the exact install command
\end{verbatim}

The design principle behind this flow is that the language model is never
the only thing standing between an error message and a command that will
later run inside a container. Four separate layers each have exactly one
job, and each layer only trusts what the previous layer has already
verified.

\subsection{Layer 1: Deterministic Retrieval}
\label{sec:m-repair-retrieval}

Before any network request, the failing Python import name must be resolved
to an installable PyPI distribution name. The two are frequently different:
\texttt{sklearn} installs as \texttt{scikit-learn}, \texttt{cv2} as
\texttt{opencv-python}, \texttt{skimage} as \texttt{scikit-image},
\texttt{umap} as \texttt{umap-learn}, \texttt{pkg\_resources} as
\texttt{setuptools}, and \texttt{Bio} as \texttt{biopython}. A small,
explicit, hand-curated mapping table (\texttt{config/package\_mapping.yaml})
holds every case the retriever is allowed to resolve; an import name absent
from this table is never assumed to also be its own package name, and the
retriever instead reports a \texttt{mapping\_unknown} status without making
any PyPI request at all. This rule matters concretely for this thesis's own
dataset: \texttt{sklearn} and \texttt{pkg\_resources} are, respectively, the
second and first most common failing modules (\Cref{tab:pd-top-modules}), so
getting this mapping right or wrong affects a large share of the dataset,
not an edge case.

Once resolved, the distribution's release metadata is retrieved from the
official PyPI JSON Simple API
(\texttt{GET /simple/\{distribution\}/}), normalising the name first
according to PEP 503. This endpoint was chosen deliberately over the older
project JSON endpoint's \texttt{releases} field, which PyPI's own
documentation marks as deprecated for new integrations; the newer endpoint
is the one PyPI recommends going forward. The returned files are grouped by
release version, and a safe candidate set is computed by dropping
unparseable files, pre-releases, and releases where every file is yanked,
and by checking each release's declared \texttt{requires-python} against
the runtime the notebook actually executed under.

That runtime version is a fixed constant, \texttt{"3.10"}
(\texttt{config/rag\_repair.yaml}), not looked up per notebook. This is not
an approximation of convenience: the Docker pipeline builds every
repository's container from the same template, with the base image
hardcoded to \texttt{python:3.10-slim}, so every notebook this component
processes genuinely ran under the same interpreter version. The
notebook's own recorded kernel metadata is deliberately not used as a
substitute, because it reflects what the notebook's original author had
installed, which is frequently stale or missing, not what the Docker
container the notebook actually failed in provides.

The final candidate list is capped to five versions, newest first, to keep
the amount of version metadata placed in the prompt bounded and readable.

\subsection{Version Incompatibility Needs More Than PyPI Metadata}
\label{sec:m-repair-compat}

For a \texttt{wrong\_version} failure -- a "cannot import name" error --
knowing that a release exists and supports the right Python version is not
enough. PyPI's own Simple API specification (PEP 691) only exposes, per
file, its filename, URL, hashes, required Python version, size, upload
time, and yanked status; it says nothing about which functions or classes a
given release actually provides. Proving that, for example,
\texttt{scipy.integrate.cumtrapz} still exists in a specific SciPy release
requires consulting that project's own documentation, not PyPI.

This gap is closed with a small, hand-curated registry
(\texttt{config/}\allowbreak\texttt{api\_compatibility\_evidence.yaml}). Each entry records one
specific import symbol, the version range in which an official source
confirms it still exists, and a link to that source. \Cref{tab:m-compat}
lists the three patterns curated for this dataset, all covering symbols
that were removed in a major SciPy or NumPy release.

\begin{table}[h]
  \centering
  \begin{tabular}{|p{4.4cm}|p{1.6cm}|p{6.0cm}|}
    \hline
    \textbf{Symbol} & \textbf{Compatible with} & \textbf{Official source} \\
    \hline
    \texttt{scipy.integrate.\allowbreak cumtrapz} & \texttt{<1.14.0} & SciPy 1.14.0 release notes, "Expired deprecations" \\
    \hline
    \texttt{scipy.sparse.\allowbreak sputils.\allowbreak isshape} & \texttt{<1.14.0} & SciPy 1.14.0 release notes, cross-checked against SciPy's source repository \\
    \hline
    \texttt{numpy.\allowbreak VisibleDeprecationWarning} & \texttt{<2.0.0} & NumPy 2.0.0 release notes, "NumPy 2.0 Python API removals" \\
    \hline
  \end{tabular}
  \caption{Curated API-compatibility evidence used for \texttt{wrong\_version} repairs.}
  \label{tab:m-compat}
\end{table}

This registry is deliberately small and manual rather than an automated
changelog crawler. Every entry was checked by hand against one official
source and dated; a pattern absent from the registry receives no fabricated
entry and simply yields no evidence, which means the affected row can only
ever be repaired as \texttt{none}. This does not generalise automatically
to unseen incompatibility patterns in a different notebook corpus -- it is
scoped to what this thesis's own dataset actually contains, and is
presented as such rather than as a general-purpose solution.

\subsection{Layer 2: What the Language Model Actually Decides}

The model receives only the grounded candidate set from Layer~1 -- never
the raw, unfiltered PyPI response -- together with the evidence summary for
a version-incompatibility case. From this, it proposes a structured JSON
object: an \texttt{action} (\texttt{install}, \texttt{pin\_version}, or
\texttt{none}), an \texttt{install\_name}, a \texttt{version} (for
\texttt{pin\_version}, drawn only from the supplied candidates), and a short
\texttt{rationale}. This is a genuine, load-bearing decision, not a
rubber-stamp: the model chooses which action the available evidence
actually supports, rather than being handed a decision already made
elsewhere and asked only to phrase it.

The schema actually validated against
(\texttt{schemas/repair\_proposal.schema.json}) is narrower than an earlier
L4 sketch that also asked the model to echo back the import name and a
\texttt{pypi\_evidence} block. Both fields were dropped once step I4 came to
implement this component: the import name is already known before the model
is ever called, so echoing it proves nothing, and re-stating retrieval data
in the model's own words only invites drift between what was actually
retrieved and what the model claims was retrieved. The schema also
enforces, through conditional rules, that \texttt{install} never carries a
version, that \texttt{pin\_version} always carries both a name and a
version, and that \texttt{none} carries neither.

Critically, \textbf{the model never returns a command field.} An earlier
draft of the prompt had the model produce the exact command string that
would later run inside a container. That design was changed once it became
clear what it implied: it would let free-text model output reach a field
executed as a subprocess. The model now proposes only the three structured
values above; the executable command is built afterwards, deterministically,
described next.

\subsection{Layers 3 and 4: Validation and Command Construction}

Before anything the model proposed is trusted, a deterministic validator
checks it against the same evidence Layer~1 already retrieved:
\texttt{action} must be one of the three allowed values; a proposed
\texttt{install\_name} must exactly match the distribution the retriever
itself already resolved, not a name the model introduced; a proposed
\texttt{version} must be literally present in the supplied candidate list,
never a nearby or plausible-looking one; a \texttt{missing\_package}
\texttt{install} additionally requires at least one candidate marked
Python-compatible; and a \texttt{wrong\_version} \texttt{pin\_version}
additionally requires that the compatibility-evidence lookup actually
resolved. Every one of these string values is also checked against a narrow
allow-list pattern that rejects a leading dash, whitespace, or shell
metacharacters, as a second, independent safety layer beneath the exact-match
checks. If any single check fails, the entire proposal is discarded and
overridden to \texttt{action: "none"} -- nothing is "corrected" to a nearby
safe value, because a corrected value is still a guess.

Only once validation passes is the actual install command constructed, and
only then, deterministically, as a list of argument tokens --
\texttt{["python", "-m", "pip", "install", install\_name]}, or with
\texttt{"install\_name==version"} for a version pin -- never as a single
shell string. A separate, human-readable \texttt{command} field is also
stored for logging, but it exists only for a reader's convenience and is
never itself executed.

If the model's response cannot be parsed, fails schema validation, fails
grounding validation, times out, or the model is unavailable, the component
retries once with the failure reason included in a corrective prompt.
Exhausting that one retry always yields \texttt{action: "none"}: an
unresolved failure is never allowed to fall through as if it had succeeded.

\subsection{Runtime Configuration and a Timeout Correction}

RAGRepairAgent shares its model with the explanation component
(\texttt{gemma2:9b} through Ollama, temperature 0.1, top-p 0.9, at most 700
generated tokens). One runtime parameter changed after real-world testing: the
request timeout was raised from an initial 120 seconds to 300 seconds. This
was not a guess -- a real pilot run (\Cref{chap:validation}) on CPU-only
hardware recorded the same repair request timing out twice at 120 seconds
and then succeeding on the first attempt once the ceiling was raised to 300
seconds, on the same machine, with the same model. The underlying model was
healthy throughout; the earlier value was simply too tight for slow,
CPU-only generation. The retrieval client also enforces a short, bounded
wait (0.5 seconds) between real PyPI requests, and at most one bounded retry
on an HTTP 429 response, so that a batch run of many rows cannot
accidentally overload the public registry.

\subsection{The Implemented Result Shape}

The result actually persisted by
\texttt{scripts/rag\_repair\_agent.py} differs from the illustrative fix
object in \Cref{subsec:fix-object}, as already noted there. Every result
record carries, among other fields, \texttt{notebook\_execution\_id} as its
only join key back to the rest of the dataset; an \texttt{eligibility} block
recording whether the row was usable, excluded, or invalid; a nested
\texttt{retrieval\_result} block holding everything Layer~1 found;
\texttt{final\_action}, \texttt{final\_install\_name}, and
\texttt{final\_version} as the validated outcome; the constructed
\texttt{argv} list and its human-readable \texttt{command} string; and an
overall \texttt{status} of \texttt{"success"} (a grounded proposal),
\texttt{"abstained"} (a legitimate \texttt{none}, whether chosen by the
model or forced by a failed check), or \texttt{"failed"} (every attempt was
exhausted without ever producing a usable response). One field the
illustrative design object showed at the top level, \texttt{import\_name},
does not appear in the persisted record at all, because it is already
present in the nested input block copied from I2 and did not need
repeating -- a detail that becomes important in \Cref{sec:m-fixapply}, since
it is exactly why FixApplicator cannot recover a notebook's identity from
an I4 record alone.

\section{Fix Application and Notebook Re-execution (I5)}
\label{sec:m-fixapply}

FixApplicator answers the question RAGRepairAgent deliberately leaves open:
\emph{does the proposed repair actually work?} It contains no language
model call anywhere -- it is a small, fully deterministic pipeline that
applies an already-validated proposal, re-runs the notebook, and classifies
what happened.

\begin{verbatim}
validated RAGRepairAgent proposal (I4)
    -> re-join with the I2 dataset by notebook_execution_id
    -> re-validate install_name/version and rebuild the install command
    -> reconstruct a Docker environment equivalent to the original
    -> apply the fix inside that environment
    -> re-run the notebook top to bottom
    -> classify the outcome
\end{verbatim}

\subsection{Recovering Notebook Identity: the I2 Re-join}

The I4 result described above is joined to the rest of the dataset by
\texttt{notebook\_execution\_id} alone; it does not carry the repository
ID, notebook name, or repository URL forward. FixApplicator therefore looks
up each I4 record's \texttt{notebook\_execution\_id} in the I2 dataset
(\texttt{data/context-classification/dependency\_error\_contexts.jsonl}) to
recover this identity before doing anything else. If a given ID is not
found, or the matching I2 record is missing one of these fields, the
attempt is recorded as an infrastructure failure (\texttt{apply\_error},
failure stage \texttt{"join"}) rather than guessed -- the alternative, filling
in a missing field with a default, would risk silently cloning or running
the wrong repository.

\subsection{Never Trusting a Persisted Command Blindly}

Even though I4 already validated its own proposal, FixApplicator does not
execute the persisted \texttt{argv} as-is. It re-checks
\texttt{install\_name} and \texttt{version} with the same safety check I4
itself uses, and rebuilds the argument list from scratch with the same
function I4 used to build it in the first place. If the freshly rebuilt
list does not match the one that was persisted, the record is rejected
before anything runs. This is a second, independent enforcement of the same
safety rule, not a redesign of it -- the two components agreeing confirms
that no tampering or drift occurred between when I4 wrote the record and
when I5 reads it.

\subsection{Rebuilding the Docker Environment}
\label{sec:m-fixapply-docker}

\Cref{chap:architecture} originally assumed FixApplicator would be handed a
container from the repository's original run. In practice, the original
per-repository containers from the upstream Docker pipeline were never kept
on this machine -- there is simply nothing left to attach to. FixApplicator
therefore rebuilds an equivalent environment from the same recipe on every
attempt: the same base image (\texttt{python:3.10-slim}), the same
\texttt{requirements.txt}-then-\texttt{setup.py} install sequence, and the
same notebook execution command
(\texttt{jupyter nbconvert --to notebook --execute --allow-errors}) the
reference pipeline itself uses. This command, not a hand-written execution
loop, is reused deliberately: it is the exact mechanism that produced the
original dataset, so re-using it keeps the comparison fair. Only the one
notebook that actually failed is executed, not every notebook in the
repository, and it is always executed from the beginning, never resumed
from the failing cell, because a notebook keeps state across cells and a
partial re-run cannot prove a fix actually works.

Two details differ from the reference pipeline on purpose. First, the build
context is kept small and separate from the cloned repository, instead of
placing build files directly inside the clone; this avoids sending an
entire git history to the Docker daemon as build context and has no effect
on the resulting image, since the repository is still bind-mounted into the
container at run time exactly as the reference pipeline itself does.
Second, the container runs as its default root user rather than mapping to
a host user ID. This is necessary because I4's already-validated install
command has no user-mapping flag of its own, and FixApplicator must never
rewrite an argument list that has already passed validation; running the
disposable, single-use container as root lets that unmodified command
succeed without inventing a different command shape than the one I4
approved.

\subsection{Commit Pinning}

Where the upstream pipeline's own database has a recorded commit for a
repository, FixApplicator checks out that exact commit before building
anything, rather than staying on the repository's current default branch.
This matters because a public GitHub repository can change after the
original dataset was collected; without pinning, a later fix attempt could
be tested against a different version of the notebook than the one that
actually produced the original failure. Notably, the reference pipeline
that originally produced this dataset does not do this itself -- it clones
only the latest state of the default branch even though it records a
commit -- so this is a case where the implementation goes further than the
system it builds on. If the recorded commit cannot be checked out, this is
treated as a specific infrastructure failure rather than a silent fallback
to the current branch. The commit lookup itself is optional: if the
upstream database is missing or cannot be read, FixApplicator simply
proceeds without commit or baseline-setup metadata for that attempt rather
than failing it outright.

\subsection{Fix Application Safety}

The install command is only ever executed as an explicit list of arguments,
passed directly to a subprocess, never assembled into a shell string. Every
external process this component invokes -- git and Docker included -- goes
through an injectable runner as an argument list; a static check of the
source code confirms that \texttt{shell=True}, \texttt{os.system},
\texttt{eval}, and \texttt{exec} appear nowhere in it. If the fix
installation itself fails inside the container, the container script stops
immediately and the notebook is never executed -- running an unfixed
notebook would validate nothing.

\subsection{Outcome Classification}

Exactly three outcomes are possible, matching the contract already fixed in
\Cref{chap:architecture}; no fourth value is ever introduced.

\begin{table}[h]
  \centering
  \begin{tabular}{|l|p{9.3cm}|}
    \hline
    \textbf{Outcome} & \textbf{Meaning} \\
    \hline
    \texttt{fixed} & The re-executed notebook produces no error output anywhere. \\
    \hline
    \texttt{still\_failing} & The fix installed and the notebook re-ran, but an error remains -- possibly the same error as before, possibly a different one further into the notebook. \\
    \hline
    \texttt{apply\_error} & The repair could not be meaningfully tested: a dataset join failure, an invalid I4 record, a clone or checkout failure, a Docker build failure, a failed fix installation, a timeout, or a missing output notebook. \\
    \hline
  \end{tabular}
  \caption{FixApplicator outcome values.}
  \label{tab:m-outcomes}
\end{table}

A \texttt{still\_failing} result additionally records whether the new error
is the same as the original one, using a conservative check: it is only
ever marked the same when both the exception type matches \emph{and} the
original failing module name still appears in the new message. This check
is designed to under-claim sameness rather than over-claim it, so that a
result is never wrongly reported as "the fix did nothing" when it actually
resolved the original problem and simply uncovered a different one. This
distinction turns out to matter a great deal in practice, as
\Cref{chap:validation} shows.

A row whose proposal was \texttt{"none"}, or whose I4 status was not
\texttt{"success"}, is not treated as an attempt at all. It is recorded with
a separate \texttt{status: "skipped"} field rather than invented as a
fourth outcome value, keeping the three-value contract intact while still
distinguishing "nothing was attempted" from "an attempt was made and
failed."

\subsection{What Real Execution Revealed That Mocked Tests Could Not}
\label{sec:m-fixapply-bugs}

The mocked, offline test suite for this component was already fully passing
before it was ever run against a real Docker daemon. Running it for real
nonetheless surfaced three environment-specific bugs that no amount of
mocked testing could have caught, because each one depends on exact runtime
behaviour that a mock necessarily replaces with an assumption:

\begin{itemize}
  \item \textbf{A locale-dependent crash while reading container output.}
        Captured subprocess output was decoded using the host machine's
        default text encoding, which silently failed on legitimate progress
        output from pip, git, and Docker and lost the entire container log
        on the very first real attempt. Fixed by decoding output as UTF-8
        explicitly, substituting any undecodable byte rather than crashing.
  \item \textbf{A corrupted script on Windows.} Writing the generated
        container entry-point script through a normal text write silently
        converted its line endings, which broke the script's first line (its
        shebang) so that Docker could not run it at all. Fixed by writing
        that file with explicit Unix-style line endings.
  \item \textbf{Leftover work directories after cleanup.} A real
        \texttt{git clone} on Windows leaves some files read-only, and the
        directory-removal call used during cleanup silently skipped its own
        error-handling logic whenever it was told to ignore errors, so
        read-only files were never actually removed. Fixed by handling
        removal errors explicitly instead of asking the cleanup call to
        ignore them.
\end{itemize}

All three are now covered by dedicated regression tests. They are recorded
here as a methodological point, not only a bug list: a mocked test suite
can prove that a component's logic is internally consistent, but it cannot
prove that the component behaves correctly against a real operating system,
a real Docker daemon, and real subprocess output. The real pilot described
in \Cref{chap:validation} was necessary precisely because of this gap.

\subsection{Environment Isolation and What Is Out of Scope}

Every attempt uses a freshly named work directory and container, and both
are always removed afterwards regardless of whether the attempt succeeded,
failed, or timed out. A rebuilt Docker image being reused across attempts
(through Docker's own build cache) is expected and is not a sign that
isolation failed -- the guarantee that matters is a fresh clone and a fresh
container per attempt, not a fresh image layer every time.

FixApplicator has no orchestrator, no two-round feedback loop, and no
persistence to a results database or knowledge graph. It returns one
self-contained result per attempt; wiring many such attempts into an
unattended batch run, feeding a \texttt{still\_failing} result back into
RAGRepairAgent for a second attempt, and writing results into the
\texttt{repair\_attempts} table and the knowledge graph are all explicitly
future work (\Cref{chap:future-work}).

\section{Design Decisions at a Glance}
\label{sec:m-decisions}

The preceding sections explain the reasoning behind each decision in place,
close to the component it affects. \Cref{tab:m-decisions} collects the most
significant decisions in one place for reference, including several where
the final choice only emerged after an initial assumption met a concrete
problem during implementation.

\begin{table}[h]
  \centering
  \begin{tabular}{|p{3.6cm}|p{9.7cm}|}
    \hline
    \textbf{Decision} & \textbf{Reasoning} \\
    \hline
    Docker pipeline, not the conda run, as integration target &
      Its residual failures, left over after containerisation, are a
      cleaner and more targeted dependency signal (\Cref{sec:pd-why-docker}). \\
    \hline
    Scope limited to \texttt{ModuleNotFoundError} / \texttt{ImportError} &
      These are exactly what the upstream pipeline already tags as
      dependency-related; no new detection logic was needed. \\
    \hline
    pip-only repair actions for v1 &
      Keeps the first version controlled and measurable; system-level and
      source-code repairs are a separate, larger problem (\Cref{sec:pd-scope}). \\
    \hline
    Deterministic classification before any LLM call &
      A decision that gates automated repair should be auditable and
      reproducible, not dependent on a model's judgement (\Cref{sec:m-dataset}). \\
    \hline
    Local, open model (Gemma-2 9B via Ollama) &
      Matches PLLM's own best configuration and keeps the system
      cost-free and privacy-preserving to run (\Cref{sec:m-explanation}). \\
    \hline
    Explanation output changed from prose to strict JSON &
      \emph{Initial:} free-form prose (L4). \emph{Problem:} an automated
      pipeline needs machine-checkable, loggable output. \emph{Final:}
      schema-validated JSON with one retry (\Cref{sec:m-explanation-deviation}). \\
    \hline
    Repair command never written by the model &
      \emph{Initial:} the model returned a command string directly.
      \emph{Problem:} this lets free-text output reach a field later run as
      a subprocess. \emph{Final:} the model proposes only structured values;
      the command is built deterministically after validation. \\
    \hline
    PyPI JSON Simple API, not the project JSON \texttt{releases} field &
      PyPI's own documentation marks \texttt{releases} as deprecated for
      new integrations (\Cref{sec:m-repair-retrieval}). \\
    \hline
    Explicit import-to-distribution mapping table &
      An import name is not always its own PyPI package name; guessing
      identity would misresolve common cases such as \texttt{sklearn}. \\
    \hline
    Curated compatibility-evidence registry &
      PyPI metadata cannot prove a specific symbol exists in a release;
      letting the model guess a historical version would be unfounded
      (\Cref{sec:m-repair-compat}). \\
    \hline
    Procedural pipeline, not an autonomous agent &
      Agentic frameworks add latency and control overhead that scale
      poorly across hundreds of independent, mostly single-dependency
      failures~\cite{yang2025aprsurvey}. \\
    \hline
    Docker environment reconstructed, not reused &
      \emph{Initial assumption:} FixApplicator attaches to the original
      container. \emph{Problem discovered:} the original containers were
      never preserved. \emph{Final:} an equivalent environment is rebuilt
      from the same recipe on every attempt (\Cref{sec:m-fixapply-docker}). \\
    \hline
    Re-join by \texttt{notebook\_execution\_id} &
      \emph{Problem discovered:} I4's own result does not carry repository
      or notebook identity forward. \emph{Final:} I5 recovers it by joining
      back to the I2 dataset before doing anything else. \\
    \hline
    Repository commit pinning &
      Prevents a later repair attempt from silently testing against a
      different revision of a public repository than the one that
      originally failed. \\
    \hline
    Full top-to-bottom re-execution &
      A notebook keeps state across cells; a partial re-run from the
      failing cell cannot prove a fix works~\cite{grotov2024debug}. \\
    \hline
    Bounded retries and bounded rounds, never unlimited &
      Applies uniformly across the explanation retry, the repair retry, the
      PyPI rate-limit retry, and the planned two-round repair mode -- a
      slow or unavailable external system must degrade to a clean failure,
      never an unbounded wait. \\
    \hline
  \end{tabular}
  \caption{Major design decisions across I1--I5, with the reasoning behind each.}
  \label{tab:m-decisions}
\end{table}
